\section{Related Work}
\label{sec:related_work}
Diffusion and score-based models have become a dominant approach to generative modeling, with continuous-time formulations connecting SDEs, probability flow ODEs, and score estimation \cite{ho2020denoising, song2021scorebased, song2020generative}. Flow matching offers a deterministic alternative that directly learns the velocity field transporting a base distribution to data, and is closely related to stochastic interpolants \cite{lipman2023flow, albergo2023building}. Our work builds on this line by studying how local perturbations propagate through learned flows and by deriving closed-form drift/Jacobian expressions in Gaussian and mixture settings.

From a causality perspective, representation learning has been proposed as a way to expose mechanistic structure in high-dimensional data \cite{scholkopf2021toward}. Complementary to this, Jacobian-based analyses have been used to encourage or measure disentanglement in generative models \cite{wei2021orthogonal, khrulkov2021disentangled}. We connect these threads by analyzing Jacobian-vector products of flow matching models and using them to estimate dependency structure in both synthetic and image domains, while explicitly distinguishing conditioning from formal interventions.
